\chapter{Очередь и ожидание}\label{ch:ochered}

\textbf{Что вы узнете:} как сделать так, чтобы один дрон начинал после другого.

\section{У дронов нет общих часов}

Помните: у каждого дрона свои часы, и они начинаются с нуля. Все дроны стартуют вместе, но они не знают, что делают другие. Длина шоу — это длительность самого долгого дрона.

\begin{figure}[H]
  \centering
  % Две ленты: у каждого дрона свои часы, длина шоу по самому долгому
\begin{tikzpicture}[x=0.52cm, font=\sffamily\small]
  \def\h{0.7}
  \node[anchor=east] at (-0.2,1.55) {дрон 1};
  \foreach \a/\b/\t/\f in {0/5/calibrate/paper, 5/10/takeoff/accentsoft, 10/11/{}/paper, 11/15/code/accentsoft, 15/18/land/paper}{
    \draw[fill=\f, draw=ink!60, thick] (\a,1.2) rectangle (\b,1.2+\h);
    \node[font=\scriptsize] at ({(\a+\b)/2},{1.2+\h/2}) {\t};
  }
  \node[anchor=east] at (-0.2,0.35) {дрон 2};
  \foreach \a/\b/\t/\f in {0/5/calibrate/paper, 5/10/takeoff/accentsoft, 10/11/{}/paper, 11/17/{set\_z: ждёт}/warnfill, 17/21/code/accentsoft, 21/24/land/paper}{
    \draw[fill=\f, draw=ink!60, thick] (\a,0) rectangle (\b,\h);
    \node[font=\scriptsize] at ({(\a+\b)/2},{\h/2}) {\t};
  }
  \draw[dashed, inkgray] (18,1.0) -- (18,2.3);
  \node[anchor=south, text=inkgray, font=\scriptsize] at (18,2.3) {дрон 1 сел, ждёт остальных};
  \draw[dashed, accent, thick] (24,-0.2) -- (24,2.3);
  \node[text=accent, font=\scriptsize\bfseries, anchor=south east] at (24.3,2.3) {конец шоу};
  \draw[-{Stealth}, ink] (0,-0.55) -- (25.3,-0.55) node[right] {$t$, с};
  \foreach \t in {0,5,10,15,20,24}{\draw[ink] (\t,-0.48) -- (\t,-0.62); \node[below, font=\scriptsize] at (\t,-0.62) {\t};}
\end{tikzpicture}

  \caption{Ленты двух дронов: первому нужно 18 секунд, второму 24. Шоу длится 24 секунды. Второй ждёт (светлый отрезок), пока первый летит.}\label{fig:clocks}
\end{figure}

Поэтому «сделай это после дрона 1» нужно устроить самому: сложить секунды дрона 1 и вставить ожидание.

\section{Эстафета}

Три дрона по очереди слетают вверх по своему столбцу и назад. Пока летит один, остальные ждут и горят синим.

\solfile{l12_relay.py}

Посчитаем время второго дрона.
\begin{itemize}
  \item К секунде 11 он закончил калибровку (5), взлёт (5) и выход на свою клетку (1).
  \item Дрон 1 с 11-й секунды летит $4 + 4 = 8$ секунд, значит освободится на 19-й.
  \item Дрон 2 ждёт \kod{set_z(100, 9)}, то есть до 20-й секунды. Секунда лишняя, это запас.
  \item Потом он летит 8 секунд, до 28-й, и дрон 3 ждёт \kod{set_z(100, 18)}, до 29-й.
\end{itemize}

\begin{zametka}
Формула ожидания: \emph{ожидание = сумма времён шагов того, кого ждёшь, + запас 1 секунда}. Если потом поменять время у первого дрона, ожидание придётся пересчитать вручную.
\end{zametka}

\begin{vnimanie}
Настоящие дроны стартуют не в один миг и летят не точно по расписанию. Оставляйте запас в секунду-две.
\end{vnimanie}

\begin{probuy}
  \item Два дрона по очереди: сначала дрон 1 летит из клетки 13 в клетку 49 и обратно (по 4\,с), потом дрон 2 из 15 в 51 и обратно.
  \item Три дрона зажигаются красным по очереди на две секунды: сначала первый, потом второй, потом третий. Остальное время они не светят.
  \item «Волна»: три дрона на 80\cm по очереди поднимаются до 150\cm и опускаются обратно (по 2 секунды вверх и вниз).
  \item Дрон 1 летит из 13 в 17 за 6 секунд, дрон 2 из 37 в 41 за 6 секунд, но на 3 секунды позже. Как записать «позже»?
\end{probuy}
\otvety{12}
